\section{Commercial discovery: compare the retained graph before the interface}
\label{exec13:commercial}

Commercial Bayesian-network systems already address substantial parts of the
workflow that a new causal product would have to build. The relevant comparison
includes interactive model construction, structure and parameter learning,
background restrictions, model exchange and programmatic inference. The
proposal should specify which of these services it would reuse and which
scientific contract it would add. A common command-line or MCP interface
alone would not establish an advantage over an existing modeling environment.

The following comparison uses inspected vendor technical documentation.
The GeNIe manual declares version 5.0.R2, built on 23 November 2024; the
SMILE programmer's manual declares version 2.2.4.R1, built on 27 April 2024.
The public URLs were retrieved after the 7 October 2026 literature cutoff,
so those headers identify the declared manuals without proving their exact
historical bytes at the cutoff. The inspected BayesiaLab pages are mutable
and do not consistently pin the technical discussion to a numbered release.
Their findings are consequently dated documentation observations. No vendor
executable, commercial trial or paid service was run for this comparison
\citep{exec13_genie_readme,exec13_smile_intro,exec13_bayesia_structure}.

GeNIe exposes both a PC route and a Bayesian Search route. The PC manual
documents a default significance level of 0.05, maximum adjacency of eight,
and no time limit by default. It allows continuous data under a multivariate
Gaussian assumption and returns a pattern for further inspection. The
manual's suggested histogram and scatterplot checks are useful exploratory
diagnostics; they cannot establish the full joint distribution or the causal
assumptions. Its Bayesian Search discussion instead describes score-guided
hill climbing with restarts, a BDeu-related sample-size parameter and a
structural prior. A seed of zero uses the processor clock, so it is not a
fixed reproducibility seed. These are different procedures to configure
and compare, even though they share a graphical application
\citep{exec13_genie_pc,exec13_genie_bs}.

The pattern editor is a particularly important capability. Its documentation
explicitly discusses an equivalence class and permits inspection of unresolved
arcs. It also documents a \emph{Create DAG} command that randomly chooses
directions for unresolved arcs while producing an acyclic graph. That command
can be useful for constructing a representative Bayesian network, but the
choice does not add observational evidence for the selected directions.
The proposed product must preserve the pre-conversion pattern and label any
selected representative. It should not describe GeNIe as necessarily forcing
a single graph: the documented pattern workflow is itself a useful alternative
\citep{exec13_genie_pattern}.

Consider the funding law in Section~\ref{exec13:survey}. Its three Gaussian
models share the same covariance but imply effects of either zero or 0.48
thousand currency units per percentage point. A proposed interface that keeps
only a randomly chosen representative could report either zero or 120 currency
units per customer for the 25 basis point action. The difference would come
from the extra orientation choice. This is a project counterexample to a
lossy integration design; it is not an observed GeNIe estimate or a claim
that the vendor's documented interface conceals that choice. A matched trial
should inspect the pattern, the conversion and the returned effect separately
\citep{exec13_examples}.

BayesiaLab supplies another relevant alternative. Its structural-learning
overview describes heuristic search and an MDL score used to compare learned
networks. The EQ page explicitly describes search over equivalence classes.
That makes class-aware search a documented competitor capability, rather
than a novelty the proposed library can claim simply by retaining ambiguous
graphs. A score optimized over graph structures still does not validate
the banking assignment mechanism or supply an intervention-effect confidence
region. Those are additional scientific questions
\citep{exec13_bayesia_structure,exec13_bayesia_eq}.

The EQ page also documents a consequential distinction among restrictions:
fixed arcs are treated as ordinary arcs, whereas forbidden arcs and temporal
indices are taken into account. A workflow that marks a bank-supplied edge
as fixed must therefore inspect the selected algorithm's actual meaning for
that setting. A generic \texttt{required\_edges} field cannot be mapped to a
similarly named graphical control without a behavioral check. The same page's
description of efficiency is a vendor claim, not a measured runtime result
or a proof of finding the global optimum in this proposal
\citep{exec13_bayesia_eq}.

\section{What a commercial integration must retain}

SMILE is the proprietary C++ engine underlying GeNIe. Its programmer's manual
describes compiled platform distributions and wrappers for Java, Python, R,
.NET and COM. These are access paths to an engine, not independent discoveries
of five competing causal implementations. Their individual interoperability
still needs review. The C++ manual does not, by itself, verify that every
wrapper exposes every endpoint type and error state
\citep{exec13_smile_intro}.

The \texttt{DSL\_pc::Learn} reference returns a \texttt{DSL\_pattern} and a
success or error code. A progress callback can interrupt learning, in which
case the documented return is \texttt{DSL\_INTERRUPTED}. The time limit is
checked after an iteration, so the documented setting is not evidence of a
strict external wall-time guarantee. A resource-bounded wrapper should retain
interruption as an incomplete computation and enforce its own lifecycle
contract if a hard deadline is required. Stopping the process cannot make
the unfinished graph a completed scientific answer
\citep{exec13_smile_pc}.

Graph conversion also requires care. The \texttt{DSL\_pattern} reference
declares absent, undirected and directed edge types, indexed by the dataset's
variable order. Its conversion to a network assumes that the dataset and
pattern coordinates correspond. GeNIe's graphical pattern documentation
uses a richer visual description, including colored and double-headed arcs.
These pages alone do not establish a complete equivalence between every
graphical mark and the C++ representation, or a standard CPDAG/PAG
serialization. Before reuse, the adapter needs named-node fixtures that
distinguish a collider, an unresolved path, a deliberately transposed
matrix and a permuted dataset. It should preserve any unresolved mapping
as an unsupported conversion
\citep{exec13_smile_pattern,exec13_genie_pattern}.

Bayesia's programmatic offering needs a different distinction. Its API overview
separates development and production configurations, and it describes a
stateless REST service for inference with an existing Bayesian network.
That description does not establish that the same REST service performs
structural discovery. The overview also describes Java configurations and
production licensing by the number of simultaneously loaded model instances.
Its table distinguishes customer-hosted Java models from the described REST
service's hosted model location. An integration proposal must select the
actual interface and operation before estimating deployment cost or data
movement. A label such as ``Bayesia API'' is too broad for either conclusion
\citep{exec13_bayesia_api}.

\section{Capabilities to reuse and capabilities the proposal still lacks}

The comparison has consequences for the proposed investment. A practical
prototype should consider an existing graphical modeling environment for
analyst exploration, then evaluate a narrow engine boundary for reproducible
requests. Rebuilding a full editor, data interface and inference engine would
need a separate justification. Table~\ref{exec13:commercial-comparison}
distinguishes documented vendor functionality from the proposal's current
obligations; it is not a numerical product ranking.

\begin{table}[htbp]
\centering\small
\caption{Documented alternatives and the proposed product's remaining work.}
\label{exec13:commercial-comparison}
\begin{tabularx}{\linewidth}{@{}P{.23\linewidth}XX@{}}
\toprule
Comparison & Inspected vendor evidence & Proposed product obligation \\
\midrule
Structure exploration & GeNIe pattern editor; BayesiaLab EQ class search &
Preserve graph class, background restrictions and conversion provenance. \\
Analyst environment & GeNIe graphical model workflow; BayesiaLab Professional desktop &
Select an existing editor or justify and build an alternative; no implemented full editor claimed. \\
Programmatic operation & SMILE C++ and named wrappers; distinct Bayesia Java/REST paths &
Select one supported path and verify target, coordinates, errors and resource behavior. \\
Uncertainty and money & Documented graph and inference objects under their model &
Retain causal identification, sampling uncertainty and monetary units as distinct requirements. \\
Deployment and support & Proprietary distributions and edition-specific licensing &
Resolve commercial rights, maintenance and full workflow cost before runtime admission. \\
\bottomrule
\end{tabularx}
\end{table}

The first three rows identify functionality the proposal should take seriously
as existing alternatives. The fourth states a proposed scientific contract;
it does not assert that a vendor lacks every possible uncertainty or causal
analysis feature. The inspected sources are too narrow for such a product-wide
absence claim. The last row concerns an additional operating obligation that
a research prototype does not already satisfy. Documentation of a commercial
offering is also not an independently verified service-level commitment
\citep{exec13_genie_pattern,exec13_smile_intro,exec13_bayesia_eq,
exec13_bayesia_professional}.

Stata provides an adjacent comparison rather than an automatic substitute for
graph discovery. Its inspected SEM introduction describes supplied path
diagrams or command-language models, linear \texttt{sem}, generalized
\texttt{gsem}, and continuous or categorical latent-variable models under
stated restrictions. These can matter for fitting a scientifically supplied
banking model. The source does not establish automated recovery of a causal
equivalence class, and it does not establish the absence of every such method
elsewhere in Stata or its community ecosystem. A community estimator using a
commercial runtime must be recorded separately from a vendor-supplied feature
\citep{exec13_stata_sem}.

\section{License, version and trial questions that remain open}

SMILE's licensing documentation separates development from deployment.
Academic use is restricted to the described academic research and teaching
context; other development use requires a commercial license. Server and
end-user deployment are separate licensing categories. A free evaluation
download therefore cannot be treated as authorization for embedding and
redistributing the engine. The source is proprietary, so the immediate reuse
candidate is a licensed engine boundary; the inspected material does not
provide an open-source implementation to port
\citep{exec13_smile_license}.

BayesiaLab Professional's licensing page describes a locally installed desktop
application, purchase or rental, and single-user or floating-token choices.
This is consistent with a separately described hosted inference API: desktop
installation and web-service deployment are different products. The API
overview describes online license validation and simultaneous-instance limits.
Exact edition, workload, support and deployment terms need to accompany a
cost comparison. No negotiated price, infrastructure cost or saving has been
measured here, and the proposal has not accepted commercial terms
\citep{exec13_bayesia_professional,exec13_bayesia_api}.

The present evidence supports a concrete next experiment: compare one
native open-source PC path with the documented GeNIe/SMILE pattern path,
using the same generated iid law and supplied background restrictions.
Preserve unresolved effects, coordinate changes, interruption and failed
requests. Evaluate BayesiaLab EQ separately as a score-based alternative,
with its own restrictions and selected interface. The output would answer
whether each chosen route preserves the agreed service and what complete
workflow effort it requires. It would not rank all vendor methods by one
point-estimate error.

The landscape remains incomplete. HUGIN technical material was not recovered
in the bounded acquisition, the attempted SAS technical locations were
unavailable, and the MathWorks search request was denied. Those are source
gaps, not findings of missing product capabilities. Further commercial
platforms, dated editions and their exact methods require their own primary
review. Public documentation can advance these comparisons now; a paid
trial is justified only after a particular unresolved behavior or cost could
change the proposed choice. The evidence supports a restricted comparison
of the documented alternatives and further source review before any wider claim.
\label{exec13:commercial-end}
